\documentclass{article}

\textwidth=6.5in
\textheight=8.5in
\voffset=-54pt
\oddsidemargin=0in
\evensidemargin=0in
\setlength{\parskip}{8pt}
\setlength{\parindent}{0pt}

\usepackage{amsmath, amssymb}
\usepackage{ifthen}

\usepackage[inline]{enumitem}
\setlist{topsep=0pt, parsep=8pt}

\usepackage[rgb,table]{xcolor}\convertcolorsUfalse
\usepackage{graphicx}

% change hyperref options
\PassOptionsToPackage{hyphens}{url}
\usepackage[bookmarks,colorlinks,linkcolor=black,citecolor=blue,pdfstartview=FitH,urlcolor=blue]{hyperref}
\hypersetup{pdfpagemode=UseNone}
\usepackage{bookmark}

\usepackage{graphicx}
\usepackage{multicol}
\usepackage{framed}
\usepackage{array}

\usepackage{icomma}

\usepackage{marginnote}
\usepackage[colorinlistoftodos]{todonotes}
\renewcommand{\marginpar}{\marginnote}

\usepackage{soul}
\usepackage[normalem]{ulem}

\usepackage{pgfplots}
\pgfplotsset{compat=1.18}  
\pgfplotscreateplotcyclelist{mycolor}{black, blue, red, green!70!black, magenta, purple, cyan, orange }  
\pgfplotsset{
  disabledatascaling,
  cycle list=tikzcolor, 
  axis lines=middle,
  every axis plot/.style={no marks, thick},
  every axis/.style={
    enlarge x limits=false,
    enlarge y limits=auto,
    xlabel style={at={(current axis.right of origin)},anchor=west},
    ylabel style={at={(current axis.above origin)},anchor=south},
    % extra x and y ticks for asymptotes
    extra x tick style={grid=major, grid style={dashed,black}},
    extra y tick style={grid=major, grid style={dashed,black}},
    /pgf/number format/1000 sep={},
  },    
  tick label style = {font=\footnotesize, fill=\currentbgcolor, inner sep=0pt},    
  xticklabel shift=1pt,    
  scaled ticks=false,
  scaled y ticks=false,
  unbounded coords=jump,
  samples=101,
  minor tick length=0,
  scatter/classes={
    open={mark=*,draw=black,fill=white, mark size=2, thin},
    closed={mark=*,draw=black,fill=black, mark size=2, thin}
  },
  width=3in,
}
\pgfplotsset{open/.style={only marks, mark=*, fill=white, thin, mark size=2}}
\pgfplotsset{closed/.style={only marks, mark=*, draw=black, fill=black,
    thin, mark size=2}}
\pgfplotsset{labeled/.style={nodes near coords, only marks, mark=*, draw=black, fill=black, thin, mark size=2, point meta=explicit symbolic}}

\usepackage{tikz}
\usetikzlibrary{
  matrix,%
  % enables the snake paths
  decorations.pathmorphing,%
  % enables bigger arrows
  decorations.markings,%
  decorations.pathreplacing,%
  decorations.text,%
  calc,%
  patterns,%
  shapes.geometric,%
  arrows,
  decorations.shapes,
  positioning,
  plotmarks,
  intersections
}
\tikzset{>=latex} % sets default arrow type in tikz
\tikzset{font=\normalsize}
\tikzset{mark size=1.5pt, mark options=thin}
\tikzset{pin distance=4pt,
  every pin edge/.style={<-, thin, shorten <= -2pt}}

\interfootnotelinepenalty=10000
\renewcommand{\thefootnote}{(\arabic{footnote})}

\usepackage{multirow, bigdelim}

% change to \usepackage[sols]{handout} to see solutions
\usepackage[sols, grading, workshop]{handout}

\newcommand\iscurrentenumi[1]{%
  \ifthenelse{\equal{#1}{\theenumi}}{}{#1}}
\setenumerate[1]{ref=\#\arabic*}
\setenumerate[2]{ref=\protect\iscurrentenumi{\theenumi}(\alph*)}
\newcommand\enumiiprefix[2]{%
  \ifthenelse{{\equal{#1}{\theenumi}}\AND{\equal{#2}{(\alph{enumii})}}}{}{}%
  \ifthenelse{{\equal{#1}{\theenumi}}\AND{\NOT{\equal{#2}{(\alph{enumii})}}}}{#2}{}%
  \ifthenelse{\equal{#1}{\theenumi}}{}{#1#2}%
}
\setenumerate[3]{ref=\protect\enumiiprefix{\theenumi}{(\alph{enumii})}\roman*}

\makeatletter

% underline links               
\let\@href=\href
\renewcommand{\href}[2]{\@href{#1}{\ul{#2}}}

\makeatother

\definecolor{purple}{rgb}{0.5,0,1.0}

\newcommand{\doedfinity}[2][\newpage]{
  \begin{problem}[#1]
    Do the Edfinity assignment ``Problem Set #2''.

    We encourage you to write down any scratch work here, as well as
    things you want your future self to remember. Nothing you write
    here will be graded; it's just for you!
  \end{problem}
}

\newcommand{\psrnewpage}{\newpage}
\newcommand{\psreflection}[2][]{

  \ifthenelse{\not\equal{#1}{nocite}}{
    \begin{enumerate}[resume]
    \item
      \begin{problem}[1in]
        Please cite any help you got on this problem set:
        \begin{itemize}[nosep]
        \item If you worked with other students, please list their names here.
        \item If you used resources other than those on our Canvas site, please list them here.
        \end{itemize}
      \end{problem}

    \end{enumerate}
  }

  \probonly{%
    #2

    \psrnewpage

    \emph{Reflection space.} It's a good habit to take a few minutes at the end of each pset to reflect on what you learned and jot down notes to your future self. Here are a few reflection questions to get you started. Nothing you write here will be graded; it's just for you!
    \begin{itemize}
    \item Take a few minutes to look back at the problems. How could you explain to someone else how to do each problem? What was your strategy, and how did you know to use that strategy? If there are problems where you don't feel able to answer this, write down which ones they are. Come back to them in a few days when we post the homework solutions!

      \vspace{1.5in}

    \item Take a look back at the learning objectives on today's worksheet; how are you feeling about each one? If there are ones you know you need more practice with, write those down here.

      \vspace{1.5in}

    \item What did you find most challenging about this material? Do you have any lingering questions about it? If so, write them down here so you don't forget them. At some point, stop by office hours or MQC to ask!

      \vfill
    \end{itemize}      
  }
}

\usepackage{fancyhdr}
\lhead{\textcolor{white}{This content is protected and may not be shared, uploaded, or distributed.}}
\rhead{}
\rfoot{\vspace*{\fill}\footnotesize\textcopyright{} \emph{Harvard Math Ma}}
\renewcommand{\headrulewidth}{0pt}
\pagestyle{fancy}
\begin{document}

\begin{center}
  \large\textsc{Problem Set 17 - 3 Differentiation Rules}
\end{center}

\begin{enumerate}
\item  \note{
    \begin{itemize}
    \item Given two polynomials $f, g$, calculate $f', g', fg, (fg)'$ and see that $(fg)' \neq f' g'$.
    \item Given a polynomial $f$, calculate $f'$, $f^2$, and $(f^2)'$ to see that $(f^2)' \neq 2f$.
    \item More derivatives using today's rules, including 2 where they are told to first simplify the function using exponent rules. 
    \end{itemize}
  }

  \doedfinity{17}

\item  

  In each part, you're given a function. Differentiate it using the Power Rule, Sum Rule, and Constant Multiple Rule. \emph{If you know other differentiation rules, you should \uline{not} use them!}

  You will need to first simplify the function; please be sure to show how you simplified it before differentiating.

  \begin{grading}
    2 points per part: 1 for simplifying and 1 for differentiating 
  \end{grading}

  \begin{enumerate}

  \item
    \begin{problem}[\vfill]
      $f(x) =\displaystyle \sqrt x \, \left( \frac{6x^2 + 7}{x}\right)
      $
    \end{problem}

    \begin{solution}
      We first use algebra to rewrite $f(x)$ as $6x^{3/2} + 7x^{-1/2}$. 
      Then we can use the power and sum rules: $f'(x) = 6(\frac{3}{2}x^{1/2}) + 
      7(-\frac{1}{2}x^{-3/2}) =$ \fbox{$9x^{1/2} - \frac{7}{2}x^{-3/2}$}
    \end{solution}

  \item

    \begin{problem}[\vfill]
      $f(x) =\displaystyle \frac{x^2 + 7}{(2x)^3}$
    \end{problem}

    \begin{solution}
      First we use algebra to rewrite $f(x)$ as $\frac{1}{8}x^{-1} + \frac{7}{8}x^{-3}$.
      Then we can use the derivative rules we've learned: 
      $f'(x) = \frac{1}{8}(-x^{-2}) + 
      \frac{7}{8}(-3x^{-4}) =$ \fbox{$-\frac{1}{8}x^{-2} - \frac{21}{8}x^{-4}$}
    \end{solution}

  \item
    \begin{problem}[\vfill]
      $f(x) = \dfrac{1}{\sqrt[4]{16x^3}} + \sqrt[5]{\pi} + \dfrac{1}{17}$
    \end{problem}

    \begin{solution}
      First we use algebra to rewrite $f(x)$ as $\frac{1}{2}x^{-3/4} + \sqrt[5]{\pi} + \tfrac{1}{17}$. Remember that $\sqrt[5]{\pi}$ and $\tfrac{1}{17}$ are constants, not
      variables. Using derivative rules, we find that $f'(x) = \frac{1}{2}(-\frac{3}{4}
      x^{-7/4})=$ \fbox{ $-\frac{3}{8} {x^{-7/4}}$}
    \end{solution}

  \end{enumerate}

\item  Let $f(x) = x^{1/3}$.

  \begin{grading}
    1 point per part
  \end{grading}

  \begin{enumerate}
  \item

    \begin{problem}[\vfill]
      Use \href{https://www.desmos.com/calculator}{Desmos} to graph $f$, and zoom in around $x = 0$. You should see that, as you zoom in, the graph looks more and more like a line. What's the slope of that line?
    \end{problem}

    \begin{solution}
      The graph of $y=x^{1/3}$ looks like this around $x=0$.

      \begin{tikzpicture}
        \begin{axis}[axis equal image, xmin=-0.3, xmax=0.3, ymin=-0.3, ymax=0.3]
          \addplot[red, domain=-0.1:0.1] {(x/abs(x))*(abs(x))^(1/3)};
        \end{axis}
      \end{tikzpicture}

      The tangent line at $x=0$ appears vertical. Vertical lines have undefined slope.
    \end{solution}

  \item
    \begin{problem}[\vfill\newpage]
      Use the Power Rule to calculate $f'(0)$. Explain why this is consistent with what you saw in (a). 
    \end{problem}

    \begin{solution}
      Using the Power Rule, $f'(x) = \frac{1}{3}x^{-2/3} = \dfrac{1}{3x^{2/3}}$, so
      $f'(0)$ is undefined. This is consistent with what we saw in (a), since the tangent
      line at $x=0$ is vertical, so the slope of the tangent line is undefined.
    \end{solution}

  \end{enumerate}

\end{enumerate}
\begin{problemonly}
  In the previous problem, since $f'(0)$ does not exist, $f$ isn't differentiable. Combining this with what we already know, we see that there are three reasons that a function $f$ could be non-differentiable at a point $a$:

  \begin{itemize}
  \item $f$ isn't continuous at $a$.
  \item The graph of $f$ has a ``corner'' at $x = a$. (That is, $f$ isn't locally linear at $a$.)
  \item The graph of $f$ is locally linear at $x = a$, but the tangent line there is vertical.
  \end{itemize}

  \begin{center}
    \rule{5in}{0.4pt}
  \end{center}
\end{problemonly}

\begin{enumerate}[resume]
\item \label{graph-(x^2-8x^2+16)}

  Let $f(x) = x^4 - 8x^2 + 16$. In this problem, you'll analyze $f$ and then sketch its graph. 

  \begin{enumerate}

  \item \label{graph-(x^2-8x^2+16)-domain}

    \begin{grading}
      0.5 points
    \end{grading}

    \begin{problem}[0.5in]
      What's the domain of $f$?  
    \end{problem}

    \begin{solution}
      All real numbers, or $(-\infty, \infty)$. 
    \end{solution}

  \item

    \begin{grading}
      0.5 points
    \end{grading}

    \begin{problem}[0.5in]
      Find the $y$-intercept of $f$.
    \end{problem}

    \begin{solution}
      $f(0) = 16$
    \end{solution}

  \item

    \begin{grading}
      1.5 points
    \end{grading}

    \begin{problem}[\stretch{0.75}]
      Find the $x$-intercepts of $f$. (Hint: Start by factoring $f(x)$ as $(x^2 - \text{something})(x^2 - \text{something})$.) 
    \end{problem}

    \begin{solution}
      \begin{align*}
        0 &= x^4-8x^2 + 16 \\
        &= (x^2 - 4)(x^2-4) \\
        &= (x+2)(x-2)(x+2)(x-2) \\
        &= (x+2)^2(x-2)^2 \\
        x &= \pm 2
      \end{align*}
    \end{solution}

  \item

    \begin{grading}
      2.5 points:
      \begin{itemize}[nosep]
      \item 1 for knowing they need to calculate $f(-x)$. (Of course, it's fine if they use a variable name other than $x$. But they shouldn't be plugging in specific values like $f(-2)$.) 
      \item 0.5 for correctly simplifying $f(-x)$
      \item 1 for a correct conclusion from their calculation of $f(-x)$. (For example, if they incorrectly end up with $f(-x) = - f(x)$ and say that $f$ is odd, they get this point.) 
      \end{itemize}
      If they just plug in a specific pair of values (like $f(-2)$ and $f(2)$) and then make a conclusion that's consistent with that, they get 1 point out of 2.5.
    \end{grading}

    \begin{problem}[\vfill\newpage]
      Algebraically determine whether $f$ is even, odd, or neither. Show your work. 
    \end{problem}

    \begin{solution}
      To see if $f$ has useful symmetry, we look at $f(-x)$:
      \begin{align*}
        f(-x) &= (-x)^4 - 8 (-x)^2 + 16 \\
        &= x^4 - 8 x^2 + 16 \\
        &= f(x)
      \end{align*}
      So, $f$ is an even function.
    \end{solution}

  \end{enumerate}

  You should have seen that $f(x) = x^4 - 8x^2 + 16$ has some useful symmetry. So, if we want to graph it, we can first focus just on graphing it on $[0, \infty)$, and then use symmetry to fill in the other half of the graph. That's what you'll do in the rest of this problem.

  \begin{enumerate}[resume]
  \item

    \begin{grading}
      1 point
    \end{grading}

    \begin{problem}[\vfill]
      Find and classify all discontinuities of $f$ on $(0, \infty)$. Justify your classifications using limits.
    \end{problem}

    \begin{solution}
      Since $f$ is given by a single (non-piecewise) formula, its discontinuities are where it's undefined. But we already saw that it's defined everywhere, so it has \fbox{no discontinuities}. 
    \end{solution}

  \item

    \begin{grading}
      1.5 points: 1 for computing $\displaystyle \lim_{x \to \infty} f(x)$ correctly, 0.5 for saying what this tells them about horizontal asymptotes 
    \end{grading}

    \begin{problem}[\vfill]
      Find all horizontal asymptotes of $f$ on $(0, \infty)$. Justify your answer by calculating all relevant limits. 
    \end{problem}

    \begin{solution}
      As $x\to \infty$, the dominant term in $x^4 - 8x^2 + 16$ is $x^4$. 
      Therefore $\displaystyle \lim_{x \to \infty} f(x) = \lim_{x \to \infty} x^4 = \infty$.
      This means that on $(0, \infty)$, $f$ has no horizontal asymptotes.
    \end{solution}

  \item

    \begin{grading}
      3 points: 0.5 for calculating $f'$ correctly, 1.5 for making a correct sign chart (or other correct reasoning), 1 for identifying the intervals correctly 
    \end{grading}

    \begin{problem}[\nstretch{1.5}]
      On what intervals in $(0, \infty)$ is $f$ increasing? On what intervals in $(0, \infty)$ is $f$ decreasing?
    \end{problem}

    \begin{solution}
      We'll make a sign chart for $f'(x) = 4x^3 - 16 x = 4x (x^2 - 4) = 4x (x - 2)(x + 2)$. $f'(x) = 0$ when $x = 0, \pm 2$. Since we're only interested in $(0, \infty)$, we can ignore $-2$ in our sign chart. By testing points, we get the following sign chart for $f'$ on $(0, \infty)$:
      \begin{center}
        \begin{tikzpicture}
          \begin{axis}[axis y line=none, x axis line style={-}, xtick={0, 2}, ymin=-2, ymax=2, height=1.25in, width=3in]
            \addplot+[domain=-2.5:5., very thin]{0};%
            \draw (axis cs: -2.5, -1.5) node[right] {sign of $f'$};%
            \draw (axis cs: -2.5, 1) node[right] {$f$};%
            \draw (axis cs: 1., -1.5) node{$-$};%
            \draw[->, xshift=2pt] (axis cs: 0.333333, 1.5) -- (axis cs: 1.83333, 0.5);%
            \draw (axis cs: 3.5, -1.5) node{$+$};%
            \draw[->, xshift=2pt] (axis cs: 2.16667, 0.5) -- (axis cs: 4.83333, 1.5);%
            \draw[->, xshift=2pt] (axis cs: 1.83333, 0.5) -- (axis cs: 2.16667, 0.5);%
          \end{axis}
        \end{tikzpicture}
      \end{center}
      This means that $f$ is increasing on $(2, \infty)$ and decreasing on $(0, 2)$. 
    \end{solution}

  \item

    \begin{grading}
      3 points, same breakdown as previous part but now for $f''$
    \end{grading}

    \begin{problem}[\vfill]
      On what intervals in $(0, \infty)$ is $f$ concave up? On what intervals  in $(0, \infty)$ is $f$ concave down?
    \end{problem}

    \begin{solution}
      We'll figure this out by making a sign chart for $f''(x) = 12 x^2 - 16 = 4 (3x^2 - 4)$.
      First, we need to determine where $f''(x)=0$.
      \begin{align*}
        4(3x^2 - 4) &= 0 \\
        3x^2 - 4 &= 0\\
        3x^2 &= 4 \\
        x^2 &= \frac{4}{3} \\
        x &= \pm \frac{2}{\sqrt{3}}
      \end{align*}
      This means that, on $(0, \infty)$, the only $x$-value for which $f'(x)$ is zero
      is $x= \frac{2}{\sqrt{3}}$.
      By testing points, we get the following sign chart for $f''$ on $(0, \infty)$:
      \begin{center}
        \begin{tikzpicture}
          \begin{axis}[axis y line=none, x axis line style={-},
            xtick={0., 1.1547}, xticklabels={$0$, $\frac{2}{\sqrt{3}}$},
            ymin=-2, height=1in, width=3in] 
            \addplot+[domain=-1.5:3., very thin]{0};
            \draw (axis cs: -1.5, -1.5) node[right] {sign of $f''$};
            \draw (axis cs: 0.57735, -1.5) node{$-$};
            \draw (axis cs: 2.07735, -1.5) node{$+$};
          \end{axis}
        \end{tikzpicture}
      \end{center}
      This means that $f$ is concave up on $(\frac{2}{\sqrt{3}}, \infty)$ and concave
      down on $(0, \frac{2}{\sqrt{3}})$.
    \end{solution}

  \item

    \begin{grading}
      1 point: 0.5 for turning point, 0.5 for inflection point. 
    \end{grading}

    \begin{problem}[\stretch{0.25}]
      What are the turning points of $f$ in $(0, \infty)$? The inflection points? 
    \end{problem}    

    \begin{solution}
      On $(0, \infty)$, $f$ has a turning point at $x=2$, since that's where $f$ changes
      from decreasing to increasing. $f$ has an inflection point at $x=\frac{2}{\sqrt{3}}$,
      since that's where $f$ changes from concave down to concave up.
    \end{solution}

  \item

    \begin{grading}
      5 points, 1 point each for the following features:
      \begin{itemize}[nosep]
      \item matches where they said $f$ should be increasing / decreasing on $(0, \infty)$
      \item matches where they said $f$ should be concave up / concave down on $(0, \infty)$
      \item $\displaystyle \lim_{x \to \infty} f(x)$ matches what they found 
      \item matches the intercepts they found
      \item has the symmetry they described
      \end{itemize}
    \end{grading}

    \begin{problem}[\vfill\newpage]
      Sketch a rough graph of $f$ on its entire domain that incorporates all of the above information. Be sure to label any significant values on the axes. 
    \end{problem}

    \begin{solution}
      Putting all of the previous parts together, we can get a reasonable graph of $f$ on $[0, \infty)$ and then use symmetry to draw the graph on $(-\infty, 0)$:
      \begin{center}
        \begin{tikzpicture}
          \begin{axis}[xtick={-2, -1.1547, 1.1547, 2}, xticklabels={$-2$,
              $-\frac{2}{\sqrt{3}}$, $\frac{2}{\sqrt{3}}$, $2$}, ytick={16}]
            \addplot+[domain=-3.5:3.5]{x^4-8*x^2+16};
            \addplot+[domain=0.4:2, red, thin]{64/3-64*x/(3*sqrt(3))};
            \addplot+[domain=-2:-0.4, red, thin]{64/3+64*x/(3*sqrt(3))};
          \end{axis}
        \end{tikzpicture}
      \end{center}
      (To make the inflection points clearer, we've drawn in the tangent
      lines at those points.)      
    \end{solution}

  \end{enumerate}

\item \label{beanstalk}

  \note{This idea will be useful for deriving the Product Rule next time.} 

  \emph{The idea of this problem will be very important next class.}

  At noon one day, Jack plants a magical bean. To his delight, it grows quickly into a giant beanstalk. Let $f(t)$ be the beanstalk's height in feet $t$ hours after noon. Suppose $f'(3) = 7$.

  \begin{enumerate}

  \item

    \begin{grading}
      2 points: 0.5 for identifying this says something about a rate at 3 pm, 1 for the units, 0.5 for saying something that makes clear this is about the rate at which the beanstalk is growing / the rate at which its height is changing 
    \end{grading}

    \begin{problem}[\vfill]
      Interpret the fact $f'(3) = 7$ in words; include units, and be sure your answer explains the meaning of both 3 and 7 in this context. 
    \end{problem}

    \begin{solution}
      At 3 pm, the beanstalk is growing at an instantaneous rate of 7 feet per hour. 
    \end{solution}

  \item \label{beanstalk-concrete-example}

    \begin{grading}
      2 points, 1 for recognizing they should multiply $f'(3)$ by the amount of time, 1 for correctly handling units (hours vs.\ minutes). 
    \end{grading}

    \begin{problem}[\vfill] 
      Roughly how much does the beanstalk grow between 3:00 pm and 3:06 pm?
    \end{problem}

    \begin{solution}
      The time between 3:00 pm and 3:06 pm is 6 minutes, or 0.1 hours. 
      Since the beanstalk is growing at an instantaneous rate of 7 feet per hour, we expect it to grow about (7 feet per hour)(0.1 hours) = \fbox{0.7 feet} between 3:00 pm and 3:06 pm. 
    \end{solution}

  \item

    \begin{grading}
      2 points
    \end{grading}

    \begin{problem}[\vfill] 
      More generally, if $h$ is a fairly small positive number, write an approximation for the amount that the beanstalk grows between $t$ hours after noon and $t + h$ hours after noon. Express your answer in function notation.
    \end{problem}

    \begin{solution}
      By the same reasoning as in \ref{beanstalk-concrete-example}, we expect the beanstalk to grow by about \[(f'(t) \text{ feet per hour})(h \text{ hours}) = \boxed{f'(t) h \text{ feet}}.\] 
    \end{solution} 

  \end{enumerate} 

\end{enumerate}
\psreflection{For next class, read \S 8.1 and 8.3.}
\end{document}
